\documentclass[11pt]{article}

\usepackage[margin=1in]{geometry}
\usepackage{amsmath}
\usepackage{amssymb}
\usepackage{booktabs}
\usepackage{graphicx}
\usepackage{subcaption}
\usepackage{hyperref}

\title{GPU ODE Solvers Across the Abstraction Stack}
\author{Draft manuscript}
\date{\today}

\input{data_macros}
\begin{document}
\maketitle

% ====================================================================
\begin{abstract}
Utkarsh et al.\ report that Julia's \texttt{DiffEqGPU.jl} is
57--322$\times$ faster than JAX for GPU ensemble ODE solves.
We reproduce their benchmark on the same Tesla~V100 hardware and
measure performance at four levels of the GPU abstraction stack:
general-purpose framework (Diffrax), idiomatic vectorized JAX,
library-generated GPU kernels (\texttt{DiffEqGPU.jl}), and
hand-written CUDA.
We find that approximately 94$\times$ of the reported gap is
attributable to framework overhead within the Python ecosystem
itself---not to a language difference.
\texttt{DiffEqGPU.jl} delivers near-CUDA performance through a
high-level API, occupying a valuable position in the stack between
idiomatic JAX and hand-written CUDA.
These results suggest that cross-ecosystem benchmarks should
control for abstraction level to avoid conflating framework
overhead with language capability.
\end{abstract}

% ====================================================================
\section{Introduction}

GPU-accelerated solving of large ensembles of ordinary differential
equations (ODEs) is important across scientific computing, from
Bayesian parameter sweeps to chemical kinetics.
Utkarsh et al.~\cite{utkarsh2023} present \texttt{DiffEqGPU.jl}, a
Julia library that generates specialized GPU kernels from high-level
ODE specifications, and benchmark it against Diffrax~\cite{diffrax}
(JAX), torchdiffeq (PyTorch), and MPGOS (C++).
On a Tesla~V100, they report that \texttt{DiffEqGPU.jl} is
57--322$\times$ faster than Diffrax for a nonstiff Lorenz ensemble.

This comparison, however, places tools at different levels of the
GPU abstraction stack against each other.
Diffrax is structured as a composable library: solver objects,
step-size controllers, and interpolation routines interact through
abstraction boundaries.
When XLA traces through these boundaries, each boundary becomes a
materialization point---an intermediate array that must be written
to global memory and read back by the next operation.
The result is multiple GPU kernel launches per time step.
This is not a deficiency of Diffrax's implementation; it is the
inherent cost of the abstractions that make it a general-purpose,
swappable-solver library.
Eliminating the overhead would require flattening the solver into
a single monolithic function, sacrificing the library's
generality.

\texttt{DiffEqGPU.jl} avoids this trade-off entirely: rather than
tracing through library abstractions, it generates a single
specialized GPU kernel (one thread per trajectory, all state in
registers) from a high-level ODE specification.
These are fundamentally different GPU execution strategies, and the
performance gap between them reflects that architectural
difference---not an inherent language limitation.

In this paper, we measure performance at four levels of what we
call the \emph{abstraction stack} (Table~\ref{tab:stack}),
reproducing the original benchmark on the same V100 hardware class
with the paper's exact codebase and hyperparameters.

\begin{table}[t]
\centering
\small
\begin{tabular}{@{}clp{6.2cm}@{}}
\toprule
Level & Name & GPU execution strategy \\
\midrule
1 & General-purpose framework &
Solver library traced through a compiler.
Abstraction boundaries between solver components create
materialization points that prevent kernel fusion;
each step compiles to multiple kernel launches with
intermediates in global memory.
\emph{Example: Diffrax (JAX/XLA).} \\[4pt]
2 & Idiomatic vectorized code &
User writes the full solver step as flat scalar arithmetic
(no library abstraction boundaries), so the compiler sees a
single fusible graph and compiles each step into one kernel
with intermediates in registers.
\emph{Example: JAX \texttt{vmap}.} \\[4pt]
3 & Library-generated GPU kernels &
Library generates a single specialized GPU kernel (one thread
per trajectory) from a high-level ODE specification.
\emph{Example: \texttt{DiffEqGPU.jl}.} \\[4pt]
4 & Hand-written GPU kernels &
User writes the CUDA kernel directly, with full control over
register usage and thread structure.
\emph{Example: CuPy RawKernel.} \\
\bottomrule
\end{tabular}
\caption{Four levels of the GPU ODE abstraction stack.
The original benchmark compared level~3 (Julia) against level~1
(Python).
Note that levels~1 and~3 both present high-level user-facing
APIs---the user writes the ODE, and the library does the rest.
The distinction is in what happens under the hood.}
\label{tab:stack}
\end{table}

An important clarification: this stack classifies \emph{GPU
execution strategy}, not user-facing API complexity.
Both Diffrax and \texttt{DiffEqGPU.jl} present high-level
APIs---the user specifies the ODE and calls a solve function.
The difference lies in the generated GPU code.

Consider a single Tsit5 step, which evaluates the right-hand side
six times and combines the results.
At level~1, Diffrax traces this through its solver abstractions
and \texttt{vmap} vectorizes across $N$ trajectories, producing
on the order of ten separate GPU kernel launches per step
(six right-hand-side evaluations plus stage management).
Between launches, stage intermediates ($N \times 3$ arrays per
stage, ${\sim}$576\,MB for $N = 8\text{M}$) must be written to
and read back from global memory (${\sim}$900\,GB/s on a V100).
At level~2, the user writes the same six evaluations as flat
scalar arithmetic; \texttt{vmap} vectorizes identically, but XLA
now sees a single fusible computation graph and compiles the
entire step into \emph{one} kernel, keeping intermediates in
registers (${\sim}$20\,TB/s effective bandwidth).
Both levels use \texttt{jax.vmap} for parallelism and
\texttt{lax.while\_loop} for adaptive time stepping---the loop
construct is incidental.
The performance difference comes from what the compiler can fuse
within each iteration.
We confirmed this by modifying the right-hand-side function
(replacing \texttt{jnp.stack} with \texttt{jnp.array} in Diffrax),
which produced no measurable change (22{,}080\,ms vs.\
22{,}078\,ms), ruling out Python-level overhead.

At level~3, \texttt{DiffEqGPU.jl} sidesteps the fusion problem
entirely: rather than tracing library abstractions, it generates a
single specialized GPU kernel (one thread per trajectory, all
state in registers) from a high-level ODE specification.
For nonstiff problems, this achieves near--level-4 performance
through a high-level API.
At level~4, hand-written CUDA kernels give full control over
register allocation and launch configuration, serving as the
hardware-floor reference.

For nonstiff Lorenz, levels~2--4 cluster tightly, and
\texttt{DiffEqGPU.jl}'s kernel generation proves highly effective.
For stiff problems, we introduce a Robertson benchmark using
the Rosenbrock23 method, where per-step linear algebra
dominates.
Using \texttt{DiffEqGPU.jl}'s batch GPU API
(\texttt{vectorized\_asolve}), levels~2--4 again cluster
tightly---\texttt{DiffEqGPU.jl} sits between the hand-written
CUDA kernel and JAX \texttt{vmap}, consistent with the nonstiff
results.
However, the high-level \texttt{solve} API adds $>$400$\times$
overhead, illustrating how API choice---not kernel
architecture---can dominate stiff GPU performance.

% ====================================================================
\section{Methodology}

\subsection{Hardware and infrastructure}

All GPU benchmarks ran on a Tesla~V100-SXM2-32GB, the same GPU
class used by the original paper.

\subsection{Benchmark problems}

\paragraph{Lorenz attractor (nonstiff).}
We use the paper's exact setup:
$\sigma = 10$, $\beta = 8/3$,
$\rho \in \mathrm{linspace}(0, 21, N)$ per trajectory, initial
condition $(1, 0, 0)$, $t \in [0, 1]$.
Fixed stepping: Tsit5, $\Delta t = 0.001$, 1000~steps.
Adaptive: Tsit5, $\mathrm{atol} = \mathrm{rtol} = 10^{-8}$,
float32.
Trajectory counts: $N$ up to $8{,}388{,}608 = 2^{23}$.

\paragraph{Robertson kinetics (stiff).}
Standard 3-state Robertson system with rate constant
$k_2 \in \mathrm{linspace}(0, 10^4, N)$ per trajectory,
initial condition $(1, 0, 0)$, $t \in [0, 10^5]$,
using the Rosenbrock23 method.
Tested at default tolerances
($\mathrm{rtol}=10^{-3}$, $\mathrm{atol}=10^{-6}$).
The parameter sweep matches the original paper's setup and
produces a mix of non-stiff ($k_2 \approx 0$) to maximally
stiff ($k_2 = 10^4$) trajectories.

\subsection{Implementations at each level}

\textbf{Level~1 (Diffrax):}
We re-ran the paper's Diffrax code, adding
\texttt{block\_until\_ready()} for correct GPU synchronization
(the original timing code lacked a synchronization barrier).
Version details are in Section~\ref{sec:versions}.

\textbf{Level~2 (JAX vmap):}
Solver written as scalar arithmetic on a single trajectory, then
vectorized with \texttt{jax.vmap}.
Fixed stepping uses \texttt{lax.fori\_loop} (unroll factor~5);
adaptive stepping uses \texttt{lax.while\_loop}---the same loop
construct as Diffrax.
Approximately 50~lines for nonstiff Tsit5 and 100~lines for the
stiff Rosenbrock23, with no library dependencies beyond JAX.

\textbf{Level~3 (\texttt{DiffEqGPU.jl}):}
We re-ran the paper's exact Lorenz script on our V100.
For the stiff Robertson problem, we used
\texttt{DiffEqGPU}~3.12.0's \texttt{vectorized\_asolve} batch
API, which pre-builds all problems on the host and solves them
in a single GPU kernel launch.
This requires \texttt{make\_prob\_compatible} to convert
problems to a GPU-compatible layout before transferring to
device memory.
Version details are in Section~\ref{sec:versions}.

\textbf{Level~4 (CUDA):}
Hand-written CUDA kernels via CuPy \texttt{RawKernel}, one thread
per trajectory, all state in registers.
For Robertson, the kernel uses an analytically derived Jacobian
and a specialized $3 \times 3$ linear solve in registers.

\paragraph{Stiff Robertson: method differences.}
Diffrax does not implement Rosenbrock methods; its closest stiff
solver is Kvaerno3, an implicit SDIRK method requiring Newton
iteration, so we exclude level~1 from the stiff comparison to
avoid a method-choice confound.
At all three remaining levels, the solver receives an analytic
Jacobian: at levels~2 and~4 this is hand-coded into the solver
(approximately 100~lines of JAX or CUDA), while at level~3 it is
passed via \texttt{ODEFunction}'s \texttt{jac} keyword and used by
the \texttt{vectorized\_asolve} kernel.
Levels~2 and~4 additionally exploit the Robertson system's
$3 \times 3$ structure with a direct linear solve (2$\times$2
reduction plus back-substitution).
The original paper benchmarked stiff Robertson at levels~1
and~3 (Diffrax and DiffEqGPU.jl), showing a 76--130$\times$
advantage for kernel generation---but did not include a level-2
comparison.

\subsection{Timing protocol}

All Python benchmarks use \texttt{block\_until\_ready()} (JAX) or
\texttt{Stream.null.synchronize()} (CuPy) to ensure the GPU has
completed before stopping the timer.
Julia benchmarks use \texttt{CUDA.@sync}.
We report the minimum of $\geq$10 timed runs after warmup.

\subsection{Validation}
\label{sec:validation}

Every benchmarked implementation is validated against float64
reference solutions computed with SciPy~\cite{scipy}:
\texttt{DOP853} (8th-order explicit) for Lorenz at $\rho = 21$
and \texttt{Radau} (5th-order implicit) for Robertson at
$t_f = 10^5$, both at near-machine-epsilon tolerances
($\mathrm{rtol} = 10^{-13}$, $\mathrm{atol} = 10^{-14}$).
Table~\ref{tab:validation} reports the maximum absolute error at
$t_\mathrm{final}$ across all state variables for each
implementation at benchmark tolerances.
All errors are consistent with float32 precision limits.
Figure~\ref{fig:validation} shows the reference solutions.

\input{validation_table}

\begin{figure}[t]
\centering
\includegraphics[width=\textwidth]{fig_validation.pdf}
\caption{Reference solutions for both benchmark problems.
Left: Lorenz final state at $t = 1$ as a function of the
ensemble parameter~$\rho$.
Right: Robertson final state at $t = 10^5$ as a function of the
rate constant~$k_2$.
All GPU implementations (not shown) agree with these curves to
within the errors reported in Table~\ref{tab:validation}.}
\label{fig:validation}
\end{figure}

% ====================================================================
\section{Results}
% Data sources for inline numbers and tables:
%   Lorenz (Diffrax rerun):      results/diffrax_latest.out
%   Lorenz (JAX vmap):           results/jax_vmap_{fixed,adaptive}.csv
%   Lorenz (Julia pre-regression): results/julia_paper_lorenz.out
%   Lorenz (CUDA kernel):        results/cuda_lorenz_{fixed,adaptive}.csv
%   Lorenz (Julia latest tested): results/julia_lorenz.out
%   Robertson (JAX vmap):        results/jax_robertson.csv
%   Robertson (CUDA kernel):     results/cuda_robertson.out
%   Robertson (Julia 2.5.1):     results/julia_robertson_25.out

\subsection{Nonstiff Lorenz}

Figure~\ref{fig:scaling} shows scaling behavior across the full
range of trajectory counts.
Table~\ref{tab:lorenz} presents the main comparison at
$N = 8{,}388{,}608$, with one row per abstraction level.
Figure~\ref{fig:bars} visualizes the gaps between adjacent levels.

\begin{table}[t]
\centering
\begin{tabular}{@{}clrr@{}}
\toprule
Level & Method & Fixed (ms) & Adaptive (ms) \\
\midrule
1 & Diffrax\textsuperscript{a}  & \DiffraxLorenzFixed & \DiffraxLorenzAdaptive \\
2 & JAX vmap                    & \JAXLorenzFixed      & \JAXLorenzAdaptive \\
3 & DiffEqGPU.jl\textsuperscript{b}  & \JuliaLorenzFixed      & \JuliaLorenzAdaptive \\
4 & CUDA kernel                 & \CUDALorenzFixed      & \CUDALorenzAdaptive \\
\bottomrule
\end{tabular}

\smallskip
{\footnotesize
\textsuperscript{a}Timed with \texttt{block\_until\_ready()}.
The paper's archived number (10{,}038\,ms fixed) did not include
GPU synchronization.
\textsuperscript{b}Pre-regression Lorenz result from the paper-era
\texttt{DiffEqGPU} stack (Julia~1.8.5). On the latest tested stack
(\texttt{DiffEqGPU~3.12.0} / Julia~1.11.3), fixed stepping still
reaches ${\sim}$132\,ms with \texttt{save\_everystep=false}, but
adaptive stepping regresses to ${\sim}$357\,ms; see
Section~\ref{sec:versions}.}
\caption{Nonstiff Lorenz Tsit5 at $N = 8{,}388{,}608$ on
Tesla~V100.}
\label{tab:lorenz}
\end{table}

\begin{figure}[t]
\centering
\includegraphics[width=\textwidth]{fig_scaling.pdf}
\caption{Wall-clock time vs.\ trajectory count for all four
abstraction levels.
The gap between Diffrax (level~1) and the other methods is
visible across all $N$; the remaining three levels converge as
$N$ grows.}
\label{fig:scaling}
\end{figure}

\begin{figure}[t]
\centering
\includegraphics[width=\textwidth]{fig_bars.pdf}
\caption{Performance at $N = 8{,}388{,}608$ with ratios between
adjacent levels.
For fixed stepping, most of the originally reported
57--322$\times$ gap sits between levels~1 and~2 ($94\times$),
attributable to framework overhead within the same language and
hardware.}
\label{fig:bars}
\end{figure}

The dominant gap is between levels~1 and~2: moving from Diffrax
to idiomatic JAX \texttt{vmap} reduces execution time by
$\RatioDiffraxJAXFixed\times$---same language, same hardware, same algorithm.
This framework overhead accounts for nearly all of the originally
reported 57--322$\times$ difference.
The level~2 to level~3 gap (JAX \texttt{vmap} to
\texttt{DiffEqGPU.jl}) is $\RatioJAXJuliaFixed\times$ for fixed stepping, which
represents the genuine advantage of \texttt{DiffEqGPU.jl}'s
specialized kernel generation.
For fixed stepping, the level~3 to level~4 gap
(\texttt{DiffEqGPU.jl} to hand-written CUDA) is $\RatioJuliaCUDAFixed\times$.
Both the CUDA kernel (${\sim}$83\% of the V100's 15.7~TFLOP/s
FP32 peak) and \texttt{DiffEqGPU.jl} (${\sim}$77\%) operate near
the hardware floor.

For adaptive stepping, the CUDA kernel (\CUDALorenzAdaptive\,ms) is $\RatioJuliaCUDAAdaptive\times$
faster than \texttt{DiffEqGPU.jl} (\JuliaLorenzAdaptive\,ms).
Both implementations take a comparable number of steps (88 vs.\
90 average), but the hand-written kernel has lower per-step
overhead from keeping all state in registers and exploiting
FSAL (First Same As Last) reuse of the final Runge--Kutta stage.

\subsection{Stiff Robertson}

Table~\ref{tab:robertson} presents the stiff Robertson comparison
across levels~2--4, all using the same Rosenbrock23
method~\cite{hairer_wanner} with a parameter sweep
$k_2 \in [0, 10^4]$.
Diffrax is omitted because it does not implement Rosenbrock
methods; its closest stiff solver (Kvaerno3, an implicit SDIRK
method) would introduce a method-choice confound.
Figure~\ref{fig:bars_stiff} visualizes the gaps.

\begin{table}[t]
\centering
\begin{tabular}{@{}clr@{}}
\toprule
Level & Method & Time (ms) \\
\midrule
2 & JAX vmap          & \JAXRobertson \\
3 & DiffEqGPU.jl\textsuperscript{a}      & \JuliaRobertson \\
4 & CUDA kernel       & \CUDARobertson \\
\bottomrule
\end{tabular}

\smallskip
{\footnotesize
\textsuperscript{a}DiffEqGPU~3.12.0, \texttt{vectorized\_asolve}
batch API with user-provided Jacobian;
see Section~\ref{sec:versions}.

Default tolerances ($\mathrm{rtol}=10^{-3}$,
$\mathrm{atol}=10^{-6}$), $N = 8{,}388{,}608$,
$k_2 \in \mathrm{linspace}(0, 10^4, N)$.
All rows use Rosenbrock23.}
\caption{Stiff Robertson Rosenbrock23 at $N = 8{,}388{,}608$ on
Tesla~V100.}
\label{tab:robertson}
\end{table}

\begin{figure}[t]
\centering
\includegraphics[width=0.6\textwidth]{fig_bars_stiff.pdf}
\caption{Stiff Robertson performance at $N = 8{,}388{,}608$.
CUDA (level~4) is $\RatioCUDAJuliaRobertson\times$ faster than
\texttt{DiffEqGPU.jl} (level~3), which is
$\RatioJuliaJAXRobertson\times$ faster than JAX \texttt{vmap}
(level~2).}
\label{fig:bars_stiff}
\end{figure}

\begin{figure}[t]
\centering
\includegraphics[width=0.6\textwidth]{fig_scaling_stiff.pdf}
\caption{Stiff Robertson scaling behavior.
All three levels show sublinear scaling at large $N$,
clustering within a small factor---consistent with the
nonstiff Lorenz results.}
\label{fig:scaling_stiff}
\end{figure}

The stiff results reinforce the level hierarchy seen in the
nonstiff case.
The CUDA kernel (\CUDARobertson\,ms), \texttt{DiffEqGPU.jl}
(\JuliaRobertson\,ms), and JAX \texttt{vmap}
(\JAXRobertson\,ms) cluster within a factor of
$\RatioCUDAJAXRobertson\times$, with \texttt{DiffEqGPU.jl}'s
kernel generation again placing it between hand-written CUDA
and idiomatic JAX.

A critical finding is the distinction between
\texttt{DiffEqGPU.jl}'s two APIs.
The high-level \texttt{solve} API---the documented entry
point for ensemble problems---dispatches each trajectory
individually, incurring host--device round-trips that
dominate at large $N$.
In our testing, \texttt{solve} takes ${\sim}$11{,}600\,ms
for the same problem, a $>$400$\times$ overhead over the
batch \texttt{vectorized\_asolve} API reported in
Table~\ref{tab:robertson}.
The batch API requires manual problem construction
(\texttt{make\_prob\_compatible}, explicit \texttt{cu()}
transfer) but achieves near-kernel-level performance by
solving all trajectories in a single GPU launch.

% ====================================================================
\section{Discussion}

\subsection{What the performance landscape reveals}

The 57--322$\times$ gap reported by Utkarsh et al.\ spans
levels~1 to~3 of the abstraction stack.
When we control for abstraction level, the genuine level~2 to
level~3 gap is $\RatioJAXJuliaFixed\times$ for fixed stepping.
The remaining $\RatioDiffraxJAXFixed\times$ is framework overhead \emph{within the
Python ecosystem itself}---Diffrax's solver abstractions prevent
XLA from fusing the step computation into a single kernel, as
discussed in Section~1.

A broader lesson for benchmark methodology: cross-ecosystem
performance comparisons should state which level of the
abstraction stack each tool occupies.
Reporting ratios across levels without this context conflates
framework overhead with language capability.

\subsection{\texttt{DiffEqGPU.jl}'s position}

\texttt{DiffEqGPU.jl} is a well-engineered library that provides
multi-backend support (NVIDIA, AMD, Intel, Apple GPUs),
event handling, and automatic differentiation---all through a
high-level API.
Its position at level~3---delivering \JuliaLorenzFixed\,ms where hand-written
CUDA gives \CUDALorenzFixed\,ms for fixed stepping---represents a genuine
engineering achievement.

This favorable position extends to stiff problems when using
the batch GPU API: for Robertson, \texttt{DiffEqGPU.jl}'s
\texttt{vectorized\_asolve} (\JuliaRobertson\,ms) sits between
the CUDA kernel (\CUDARobertson\,ms) and JAX \texttt{vmap}
(\JAXRobertson\,ms), consistent with the nonstiff ordering.

However, the choice of API and version is critical.
The high-level \texttt{solve} interface---which handles
ensemble dispatch, output collection, and callback
support---adds $>$400$\times$ overhead for stiff problems,
placing level~3 far behind level~2.
The batch \texttt{vectorized\_asolve} API bypasses this
dispatch layer, achieving competitive performance at the cost
of requiring manual problem construction.
Both the nonstiff Lorenz and stiff Robertson results use batch
APIs; however, Table~\ref{tab:lorenz} retains the paper-era
Lorenz numbers because the latest tested \texttt{DiffEqGPU}
stack still regresses adaptive nonstiff performance
substantially. Fixed stepping can be recovered on the latest
stack with explicit \texttt{save\_everystep=false}, but
adaptive stepping remains much slower
(Section~\ref{sec:versions}).

\subsection{Reproducibility and version sensitivity}
\label{sec:versions}

Python-side benchmarks used JAX~0.9.2 and Diffrax~0.7.2 (with
\texttt{block\_until\_ready()} for correct GPU synchronization).
The Julia results use two different \texttt{DiffEqGPU} versions,
reflecting the state of the library's GPU APIs:

\begin{itemize}
\item \textbf{Lorenz (nonstiff):} \texttt{DiffEqGPU~2.1.0} with
Julia~1.8.5, using the \texttt{vectorized\_solve}/\texttt{vectorized\_asolve}
batch API and \texttt{cu(probs)} to transfer problems to the GPU.
Table~\ref{tab:lorenz} keeps these pre-regression numbers because,
on the latest tested stack (\texttt{DiffEqGPU~3.12.0},
Julia~1.11.3, CUDA.jl~5.11.0 on Tesla~V100), adaptive nonstiff
performance remains much slower:
\texttt{vectorized\_asolve} takes ${\sim}$357\,ms instead of
54\,ms. Fixed stepping is not the limiting issue on the current
stack: with explicit \texttt{save\_everystep=false},
\texttt{vectorized\_solve} still runs in ${\sim}$132\,ms,
essentially matching the pre-regression 131\,ms result. We did
not reproduce nonstiff parameter corruption from
\texttt{make\_prob\_compatible} on this tested stack. The
high-level \texttt{solve} API remains much slower
(${\sim}$22{,}000\,ms fixed), so the reported Julia Lorenz
numbers continue to come from lower-level batch APIs.

\item \textbf{Robertson (stiff):} \texttt{DiffEqGPU~3.12.0} with
Julia~1.11.3, using the \texttt{vectorized\_asolve} batch API
with \texttt{make\_prob\_compatible}.
The paper's pinned stack (2.1.0) does not export stiff GPU-kernel
solvers; the earliest version with stiff support (2.5.1) lacks
\texttt{make\_prob\_compatible}, causing \texttt{cu(probs)} to
silently produce incorrect results.
\end{itemize}

The Python-side results are stable: informal testing across
JAX~0.4 through~0.9 and Diffrax~0.5 through~0.7 suggests timings
within ${\sim}$15\% of those reported here.
Running Julia~1.11.3 on the V100 required a workaround
for a bug in \texttt{CUDA\_Driver\_jll}~13.2, which loads a
forwards-compatible driver that drops Volta
support.\footnote{\texttt{JULIA\_CUDA\_USE\_COMPAT=false}
bypasses this.}

\subsection{Limitations}

This study is specific to the Tesla~V100 and the
Lorenz/Robertson problems.
The stiff CUDA and JAX \texttt{vmap} Rosenbrock23 implementations
exploit the Robertson system's $3 \times 3$ structure with a
direct linear solve; \texttt{DiffEqGPU.jl} receives an analytic
Jacobian but uses a general-purpose linear solver.
Generalizing the level-2 and level-4 implementations to
arbitrary systems would require auto-generated Jacobians and
adaptive linear solvers.
All GPU benchmarks use float32; float64 comparisons may
differ due to the V100's 32:1 FP64:FP32 throughput ratio.

% ====================================================================
\section{Conclusion}

The 57--322$\times$ gap reported between Julia and JAX for GPU
ensemble ODE solving reflects a comparison across abstraction
levels, not a language performance difference.
When measured at comparable positions on the stack, the genuine
advantage of \texttt{DiffEqGPU.jl}'s kernel generation over
idiomatic JAX is $\RatioJAXJuliaFixed\times$ for fixed stepping---meaningful, but
far from the headline figure.
\texttt{DiffEqGPU.jl} occupies a valuable position at level~3,
delivering near-hardware-floor performance for both nonstiff and
stiff problems through its kernel-generation approach.
For stiff problems, however, the choice of API matters: the batch
\texttt{vectorized\_asolve} API achieves competitive performance,
while the high-level \texttt{solve} API adds $>$400$\times$
overhead---a finding that highlights how dispatch-layer costs can
dominate kernel performance in ensemble GPU workloads.
More broadly, cross-ecosystem benchmarks should identify each
tool's position on the abstraction stack to avoid presenting
framework-overhead ratios as language performance gaps.

% ====================================================================
\section*{Acknowledgments}

Compute resources were provided by the VSC (Flemish
Supercomputer Center), funded by the Research Foundation --
Flanders (FWO) and the Flemish Government, on the KU~Leuven
Tier-2 Genius cluster.

% ====================================================================
\begin{thebibliography}{9}

\bibitem{utkarsh2023}
U.~Utkarsh, V.~Churavy, Y.~Ma, T.~Besard, P.~Srisuma,
T.~Gymnich, A.~R.~Gerlach, A.~Edelman, G.~Barbastathis,
R.~D.~Braatz, and C.~Rackauckas,
\emph{Automated Translation and Accelerated Solving of
Differential Equations on Multiple GPU Platforms},
Computer Methods in Applied Mechanics and Engineering,
vol.~418, 2024.
\newblock \url{https://doi.org/10.1016/j.cma.2023.116591}

\bibitem{gpuodebenchmarks_repo}
U.~Utkarsh,
\emph{GPUODEBenchmarks},
\url{https://github.com/utkarsh530/GPUODEBenchmarks}.

\bibitem{diffeqgpu_repo}
SciML,
\emph{DiffEqGPU.jl},
\url{https://github.com/SciML/DiffEqGPU.jl}.

\bibitem{diffrax}
P.~Kidger,
\emph{On Neural Differential Equations}, PhD thesis,
University of Oxford, 2022.
\newblock Diffrax: \url{https://github.com/patrick-kidger/diffrax}.

\bibitem{scipy}
P.~Virtanen et al.,
\emph{SciPy 1.0: Fundamental Algorithms for Scientific Computing
in Python},
Nature Methods, vol.~17, pp.~261--272, 2020.
\newblock \url{https://doi.org/10.1038/s41592-019-0686-2}

\bibitem{hairer_wanner}
E.~Hairer and G.~Wanner,
\emph{Solving Ordinary Differential Equations~II: Stiff and
Differential-Algebraic Problems},
Springer Series in Computational Mathematics, vol.~14,
2nd ed., Springer, 1996.

\bibitem{vsc}
VSC -- Flemish Supercomputer Center,
\url{https://www.vscentrum.be/}.

\end{thebibliography}

\end{document}
